\section{Positive tensor densities for hook diagrams}\label{sec:density}
\label{ac:part:120}

The sparse estimate settles small representation dimensions. For the
remaining diagrams, the induction will compare row and column symmetry on
a rectangle with holes. We construct positive tensor densities that
dominate the projections for hook diagrams. Their trace-one normalization
will permit that comparison even after sites have been removed.

Let \(p\in[1,s]\) be an integer and let \(V=E\oplus O\), where \(E\) and
\(O\) each have dimension \(p\). On \(V^{\otimes l}\), \(l\le s\), let
\(S_l\) permute the slots with signs: interchanging two neighboring odd
factors (vectors in \(O\)) incurs a minus sign, and the other neighboring
interchanges do not. On parity-homogeneous tensors the sign is therefore
the sign of the induced reordering of the odd factors. This defines a
unitary action. An \emph{even density} is a positive semidefinite operator
on \(V\) of trace one that is block diagonal for \(E\oplus O\).

\begin{lemma}[Signed density domination]\label{lem:density}
Let \(P_\alpha\) be the isotypic projection for \(\alpha\vdash l\) in this
signed tensor representation. The types present are precisely those in the
\((p,p)\)-hook: all their boxes lie in the first \(p\) rows or the first
\(p\) columns. For every such type there is a probability distribution
\(\mu_\alpha\) on even densities such that, in positive semidefinite order,
\begin{equation}
 P_\alpha\preceq
 \exp\{F(\alpha)+O(p^2\log(s+1))\}
 \mathbb E_{r\sim\mu_\alpha}r^{\otimes l}.
 \label{ac:eq:13}
\end{equation}
Its multiplicity is at most \(\exp(O(p^2\log(s+1)))\), and the number of
types present is at most \((s+1)^{2p}\). All constants are absolute.
\end{lemma}

The signed tensor action and its hook support are classical in the hook
theory of Berele and Regev~\cite{BereleRegev1983}. Positive domination by
mixtures of tensor powers also underlies the postselection technique of
Christandl, K\"onig and Renner~\cite{ChristandlKonigRenner2009}. We derive
the needed signed form from ordinary Schur--Weyl sectors, keeping the
irreducible dimension \(D_\alpha\) explicit.

\begin{proof}
For \(l=0\), the tensor space and its only representation have dimension
one, and any density gives the assertion. Suppose \(l>0\). Fix the number
\(t\) of even factors. A fixed parity pattern has the ordinary tensor
powers of \(E\) and \(O\), with the symmetric-group action on the latter
twisted by sign. Summing over parity patterns induces this action from
\(S_t\times S_{l-t}\) to \(S_l\). Write \(W_\sigma(E)\) and \(W_\tau(O)\)
for the ordinary Schur--Weyl multiplicity spaces, the polynomial
irreducibles of highest weights \(\sigma\) and \(\tau\). The resulting
orthogonal sectors, indexed by \(\sigma\vdash t\) and \(\tau\vdash l-t\)
with at most \(p\) rows each, are
\[
 \mathcal V_{\sigma,\tau}\cong
 \operatorname{Ind}_{S_t\times S_{l-t}}^{S_l}
       ([\sigma]\boxtimes[\tau'])
 \otimes W_\sigma(E)\otimes W_\tau(O).
\]
This decomposition records the commuting actions of \(S_l\) and
\(U(E)\times U(O)\): the induction records the positions of the two
parities, and the transpose on \(\tau\) records the odd sign twist.

The representation \([\tau']\) occurs upon inducing sign representations
on subgroups of sizes \(\tau_1,\ldots,\tau_p\), since its columns can be
added successively as vertical strips. Thus, if \(\alpha\) occurs in
\(\mathcal V_{\sigma,\tau}\), Pieri's rule reaches it from \(\sigma\) by
at most \(p\) vertical strips. Transposing and interchanging the two
factors gives the analogous statement starting from \(\tau\). Padding by
zeros,
\[
 \alpha_i\le\sigma_i+p,\qquad \alpha'_j\le\tau_j+p.
\]
In particular \(\alpha_{p+1}\le p\), the hook condition. Conversely, for
a hook shape take its first \(p\) rows as the even shape, then add its boxes
below those rows by columns from left to right. These are at most \(p\)
vertical strips. In the odd tensor power, words with one basis symbol for
each strip and with its prescribed multiplicity carry the corresponding
induction of sign representations. Pieri's rule therefore supplies the
given hook shape. The first \(p\) row lengths and first \(p\) column
lengths specify any hook shape, so there are at most \((s+1)^{2p}\) of
them.

We next compare the weight of a sector with the dimension of any type
\(\alpha\) occurring there. Put
\[
 \mathcal S=l\log l-\sum_i\sigma_i\log\sigma_i-\sum_j\tau_j\log\tau_j,
\]
with zero summands interpreted as zero. The horizontal and vertical arm
lengths of \(\alpha\) outside its \(p\)-by-\(p\) core are
\[
 H_i=(\alpha_i-p)_+\quad(i\le p),\qquad
 V_j=(\alpha'_j-p)_+\quad(j\le p).
\]
The strip inequalities give \(H_i\le\sigma_i\) and \(V_j\le\tau_j\).
Each hook in a horizontal arm is at most its row hook plus \(p\).
Consequently an arm of length \(H_i\) has hook product at most
\[
 \prod_{r=1}^{H_i}(r+p)
 =H_i!\binom{H_i+p}{p}
 \le H_i!(s+p)^p.
\]
The same bound holds for each vertical arm, and the at most \(p^2\) core
hooks have length at most \(s\). The hook-length formula now gives
\[
 D_\alpha\ge(s+p)^{-O(p^2)}
       \frac{l!}{\prod_i\sigma_i!\prod_j\tau_j!}.
\]
Using \(\log(n!)=n\log n-n+O(\log(n+1))\), whose linear terms cancel
because \(|\sigma|+|\tau|=l\), proves
\begin{equation}
 \mathcal S\le F(\alpha)+O(p^2\log(s+1)).
 \label{ac:eq:14}
\end{equation}

To obtain a positive density for the sector, define
\[
 d_E=\operatorname{diag}(\sigma_1/l,\ldots,\sigma_p/l),\qquad
 d_O=\operatorname{diag}(\tau_1/l,\ldots,\tau_p/l),
\]
and let \(r_U=(U_Ed_EU_E^*)\oplus(U_Od_OU_O^*)\), with independent Haar
unitaries on \(E\) and \(O\). This is an even density. Write
\(m_\sigma=\dim W_\sigma(E)\) and \(m_\tau=\dim W_\tau(O)\). The Weyl
dimension formula gives, for \(|\sigma|=t\),
\[
 m_\sigma
 =\prod_{1\le i<j\le p}
       \frac{\sigma_i-\sigma_j+j-i}{j-i}
 \le(t+1)^{p(p-1)/2},
\]
and the analogous bound for \(m_\tau\), with \(t\) replaced by \(l-t\).
The tensor power of \(d_E\) is positive and acts on the multiplicity
space as \(W_\sigma(d_E)\). Its highest-weight line has eigenvalue
\(\prod_i(\sigma_i/l)^{\sigma_i}\); all other eigenvalues are
nonnegative. The analogous statement holds for \(d_O\). These assertions
extend to zero eigenvalues by continuity, with \(0^0=1\).

The tensor density preserves every Schur--Weyl sector. Averaging over the
two unitary groups makes its action scalar on the two irreducible
multiplicity spaces. On the whole sector \(\mathcal V_{\sigma,\tau}\) the
scalar is
\[
 \frac{\Tr W_\sigma(d_E)}{m_\sigma}
 \frac{\Tr W_\tau(d_O)}{m_\tau}
 \ \ge\
 \frac{\prod_i(\sigma_i/l)^{\sigma_i}
       \prod_j(\tau_j/l)^{\tau_j}}{m_\sigma m_\tau}
 =\frac{e^{-\mathcal S}}{m_\sigma m_\tau}
 \ge e^{-\mathcal S-O(p^2\log(s+1))}.
\]
The same scalar acts on every parity pattern: the density preserves those
patterns and has the same tensor factors on each. Thus their number does
not enter the scalar. If \(Q_{\sigma,\tau}\) is the orthogonal projection
onto this sector, positivity and \eqref{ac:eq:14} give
\[
 Q_{\sigma,\tau}\preceq
 \exp\{F(\alpha)+O(p^2\log(s+1))\}\,\mathbb E_U r_U^{\otimes l}
\]
whenever \(\alpha\) occurs there.

The relevant sector projections cover \(P_\alpha\). Their number
\(N_\alpha\) is at most \((s+1)^{2p}\), because the two padded partitions
are specified by \(2p\) integers between zero and \(s\). Sum the last
bound over these sectors, and choose a sector uniformly before choosing
its two Haar unitaries. The resulting probability distribution
\(\mu_\alpha\) satisfies
\[
 \sum_{\text{relevant sectors}}\mathbb E_U r_U^{\otimes l}
 =N_\alpha\,\mathbb E_{r\sim\mu_\alpha}r^{\otimes l}.
\]
The factor \(N_\alpha\) is absorbed by \(O(p^2\log(s+1))\), proving
\eqref{ac:eq:13}. Finally, taking traces proves the multiplicity bound:
\(\Tr P_\alpha=D_\alpha\operatorname{mult}(\alpha)\), while every tensor
density has trace one.
\end{proof}

When sites are partitioned into blocks, signed unitary reordering makes
the action of their product subgroup a tensor product. Applying
\eqref{ac:eq:13} in each block therefore gives products of densities
constant within the blocks. Each density is even, so undoing the signed
reordering carries this to the ordinary tensor product with its sites
relabeled. In particular, a constant tensor power of an even density
commutes with the full signed symmetric-group action.
